\documentclass[11pt]{article}

\usepackage[utf8]{inputenc}
\usepackage[T1]{fontenc}
\usepackage{lmodern}
\usepackage{geometry}
\usepackage{amsmath,amssymb,amsfonts,amsthm,mathtools}
\usepackage{bm}
\usepackage{enumitem}
\usepackage{microtype}
\usepackage{hyperref}
\usepackage{titlesec}
\usepackage{booktabs}
\usepackage{array}
\usepackage{setspace}
\usepackage{mathrsfs}
\usepackage{physics}

\geometry{margin=1in}
\hypersetup{colorlinks=true, linkcolor=black, urlcolor=blue, citecolor=black}
\setlist[enumerate]{topsep=0.4em,itemsep=0.35em}
\setlist[itemize]{topsep=0.3em,itemsep=0.25em}
\titleformat{\section}{\large\bfseries}{\thesection.}{0.5em}{}
\titleformat{\subsection}{\normalsize\bfseries}{\thesubsection.}{0.5em}{}
\setstretch{1.06}

\newcommand{\Aether}{\AE{}ther}
\newcommand{\Aflow}{\AE{}ther-flow}
\newcommand{\TheoryName}{\Aflow{} Interpretation of Relativity}
\newcommand{\TheoryProgram}{\Aether{} / \Aflow{} framework}

\newtheorem{definition}{Definition}
\newtheorem{proposition}{Proposition}
\newtheorem{remark}{Remark}
\newtheorem{corollary}{Corollary}

\title{\textbf{The \TheoryName{}}\\[0.4em]
\large Flagship Public Article}
\author{}
\date{}

\begin{document}

\maketitle

\begin{abstract}
This article is a single-entry public synthesis of \TheoryName{} built on the active exact-closure manuscript sequence. Its job is not to replace that sequence, but to make it readable to first-time visitors, technically trained readers, and future collaborators. The central public message is simple. The flagship deliverable of the repository is already the overview-first exact-closure package centered on the \emph{Exact-Closure Sequence Overview}. Within that package, the effective gravitational dynamics are adopted to be exactly those of general relativity with universal matter coupling, while the \Aether{} / \Aflow{} ontology supplies the interpretive and conceptual setting. At observable scale, the benchmark package is operationally identical to ordinary GR: Einstein--Hilbert dynamics, one operative metric, universal matter coupling, ordinary diffeomorphism gauge structure, standard null / proper-time / redshift structure, and no extra preferred-frame, vector, or scalar graviton sector. It is therefore not Einstein-aether theory, not Ho\v{r}ava gravity, and not a new low-energy alternative to GR.

The article makes five points explicit. First, the active theory statement is exact relativistic closure, not a low-energy deformation of GR. Second, the \Aether{} / \Aflow{} language is retained in disciplined ontological form through the notions of underlying substrate, intrinsic ordered motion, observed three-dimensional space as local experiential slice, and S-time as experienced order of change. Third, the public claim boundary is strict: GR is adopted exactly in the active package, but a completed first-principles substrate derivation is not yet claimed. Fourth, the overview, exact-closure note, and five companion manuscripts form the ordered scientific core of the package. Fifth, the present article is a release-facing wrapper built on that core, while downstream derivational materials remain subordinate status and future-work context rather than part of the benchmark claim itself.
\end{abstract}

\section{Introduction}

The repository now has a stable public-facing theory package. That fact changes the right mode of presentation. The project no longer needs to be introduced primarily as ontology awaiting its first disciplined theory statement. The exact theory statement already exists in the overview-first exact-closure sequence. What is now needed is a public article that explains that package cleanly, states the claim boundary honestly, and helps readers know where to start.

This article therefore has a narrower role than the companion manuscripts. The \emph{Exact-Closure Sequence Overview} remains the canonical front door to the package itself, and the \emph{Exact-Closure Note} remains the shortest standalone anchor. The present article is instead a public-facing synthesis. It explains what the active package is, what it does and does not claim, how the companion papers fit together, and why the larger derivational archive should be read only afterward.

At observable scale, the package keeps the Einstein--Hilbert law, one operative metric, universal matter coupling, the ordinary ADM / diffeomorphism constraint structure, standard null / proper-time / redshift rules, and no extra preferred-frame, vector, or scalar graviton sector \cite{Einstein1916,ADM1962}. The congruence-based \Aflow{} layer belongs to standard GR observer kinematics \cite{Raychaudhuri1955,EllisVanElst1999}, not to a second observer-level field sector. For that reason the benchmark package should be classified immediately as not Einstein-aether theory, not Ho\v{r}ava gravity, and not any other low-energy preferred-structure alternative \cite{JacobsonMattingly2001,JacobsonMattingly2004,ElingJacobson2004,Jacobson2010,Jacobson2014,JacobsonSperanza2015,Horava2009}. The remaining derivational ambition belongs instead inside the established emergent, induced, reconstruction, self-coupling, and analogue-gravity conversation \cite{Jacobson1995,Visser2002,HojmanKucharTeitelboim1976,Deser2004,Lovelock1971,BarceloLiberatiVisser2011,Sakharov1967}.

\paragraph{Framework and claim boundary.}
\TheoryProgram{} denotes the broader \Aether{} / \Aflow{} framework built on the claims that the \Aether{} is the underlying four-dimensional substrate of reality, the \Aflow{} is its intrinsic ordered motion, observed three-dimensional space is the local experiential slice of that deeper substrate, S-time is the experienced order of change arising from matter, light, and the \Aflow{}, and observed expansion is the three-dimensional appearance of deeper four-dimensional ordered motion. Within that framework, \TheoryName{} denotes the exact-closure theory in which the effective gravitational dynamics are adopted to be exactly Einsteinian with universal matter coupling. In the active sequence, \emph{adoption} means use of that established relativistic dynamics without claiming substrate derivation, while \emph{derivation} is reserved for a first-principles recovery from explicit substrate variables.

The immediate public task is therefore to package, review, and present the exact-closure sequence as the achieved flagship result. Downstream derivational materials remain scientifically relevant, but they belong to the later status layer of the project and should not control the opening presentation of the benchmark package.

\paragraph{What is new here?}
What survives as novelty in the current package is interpretive and architectural, not a claim of new low-energy dynamics. The package keeps adoption sharply separated from derivation, makes the overview-first benchmark sequence the front door, formalizes the \Aflow{} vocabulary through a congruence-based interpretive dictionary, and keeps the positive benchmark result separate from the downstream bounded no-go on the frozen derivational line. None of that introduces a second observer-level gravitational law beyond GR.

\section{What This Package Is Not}

For first-time readers, the comparison boundary is deliberately blunt.
\begin{itemize}
    \item \textbf{Not Einstein-aether theory.} The benchmark package does not add a dynamical preferred-frame vector field in the low-energy action; the congruence / \Aflow{} layer is a kinematical GR dictionary inside the adopted metric spacetime.
    \item \textbf{Not Ho\v{r}ava or preferred-foliation gravity.} The active package does not claim a preferred slicing, modified low-energy symmetry structure, or an extra low-energy foliation sector.
    \item \textbf{Not a low-energy alternative to GR.} At observable scale the package is ordinary GR with one operative metric and universal matter coupling.
    \item \textbf{Not a completed derivation of GR.} Exact relativistic adoption is the benchmark result now; first-principles substrate derivation remains open downstream.
    \item \textbf{Not a still-live frozen-line program.} The current frozen primitive-reservoir line already ends in the scoped verdict \texttt{Not Derived On Current Line}, so another same-output relay below that frozen object is not the current front-line task.
\end{itemize}

The positive mirror is equally simple: the exact-GR benchmark package is the public result now, the derivational burden is downstream and still open, and the frozen current line is already preserved under \texttt{Not Derived On Current Line}.


\section{Public Role of This Article}

\begin{definition}[Flagship public article]
A flagship public article is a release-facing manuscript that introduces the active exact-closure package in one readable document, preserves the fixed claim boundary of the package, and directs readers into the full manuscript sequence without reopening the underlying theory line.
\end{definition}

\begin{proposition}[Sequence-preserving role]
The present article is secondary to the exact-closure sequence overview in scientific authority and sequence order, but primary as a public orientation document for readers who need a single coherent entry article.
\end{proposition}

This distinction matters. If the public article were allowed to replace the exact-closure overview, the package could drift into a looser or more promotional presentation. If it were omitted entirely, first-time readers would be left to infer the status of the project from a distributed manuscript tree. The correct middle position is to let the package remain the source of truth while the present article functions as its public guide.

\section{Ontology in Disciplined Form}

The conceptual core of \TheoryProgram{} is ontological rather than deformation-first. The active vocabulary is:
\begin{itemize}
    \item \Aether{}: the underlying four-dimensional substrate of reality;
    \item \Aflow{}: the intrinsic ordered motion of that substrate;
    \item observed three-dimensional space: the local experiential slice of the deeper substrate;
    \item S-time: the experienced order of change arising from matter, light, and the \Aflow{};
    \item observed expansion: the three-dimensional appearance of deeper four-dimensional ordered motion.
\end{itemize}

These terms are intended to discipline the interpretive proposal, not to overstate the dynamical derivation. The \Aflow{} is not introduced as a naive vector wind in ordinary observed space. The observed world is not identified with the whole ontology. The point of the framework is instead that the observer-accessible relativistic world may be interpreted as the experienced appearance of deeper ordered substrate structure.

\begin{remark}
At the active public stage, this ontology is real as interpretive structure, but it is not yet advertised as a completed microphysical derivation of the relativistic sector.
\end{remark}

\section{Exact Relativistic Theory Statement}

\begin{definition}[\TheoryName{}]
\TheoryName{} is the exact-closure, GR-consistent theory of \TheoryProgram{} in which the \Aether{} / \Aflow{} ontology is retained while the effective gravitational dynamics are exactly those of general relativity with universal matter coupling.
\end{definition}

Its effective action is
\begin{equation}
S_{\mathrm{eff}}
=
\frac{c^3}{16\pi G}\int d^4x\,\sqrt{-g}\,R
+
S_{\mathrm{matter}}[g,\psi],
\label{eq:Seff}
\end{equation}
with field equations
\begin{equation}
G_{\mu\nu}
=
\frac{8\pi G}{c^4}T_{\mu\nu},
\label{eq:Einstein}
\end{equation}
the standard null condition
\begin{equation}
0=g_{\mu\nu}\,dx^\mu dx^\nu,
\label{eq:null}
\end{equation}
and the standard proper-time rule
\begin{equation}
d\tau^2=-\frac{1}{c^2}g_{\mu\nu}\,dx^\mu dx^\nu.
\label{eq:tau}
\end{equation}

\begin{proposition}[Operational identity with GR]
For fixed matter sector, initial data, and boundary conditions, the effective observable content of \TheoryName{} is exactly that of general relativity.
\end{proposition}

This proposition is the center of the exact-closure package. The project does not currently claim a verified low-energy deformation of GR. It claims that the ontology can be retained while the effective gravitational law remains exactly Einsteinian.

\section{Adoption and Explicit Nonclaims}

\begin{definition}[Adoption]
Adoption is the use of an already established effective dynamics, here GR with universal matter coupling, without claiming that the dynamics have already been derived from deeper substrate variables.
\end{definition}

\begin{definition}[Derivation]
Derivation is a first-principles recovery of the effective relativistic sector from explicit substrate variables, together with whatever mode control, uniqueness, and closure arguments are needed to justify that recovery.
\end{definition}

The active public claim of \TheoryName{} is adoption with exact closure, not completed derivation. That boundary should be stated positively rather than defensively. It means that the effective theory statement is already disciplined and complete at the operational level, while the deeper foundational recovery problem remains open.

The package therefore establishes the following.
\begin{enumerate}
    \item The ontological vocabulary is fixed in disciplined form.
    \item The exact effective relativistic law is fixed as Einsteinian with universal matter coupling.
    \item The package carries no separate established low-energy non-GR observable sector.
    \item The project distinguishes clearly between interpretation, adoption, and derivation.
\end{enumerate}

The package also does \emph{not} claim the following.
\begin{enumerate}
    \item It does not claim that GR has already been derived from substrate microphysics.
    \item It does not claim that ontology by itself settles the microphysical problem.
    \item It does not claim that the larger bridge archive has already produced a completed Einstein derivation.
    \item It does not claim that the present public package must be read through the bridge archive before it can count as a theory statement.
\end{enumerate}

\section{Downstream Derivational Status}

The benchmark package is already fixed, so any downstream derivational manuscript must be read against it rather than as a replacement front door. On the current frozen primitive-reservoir line, the bounded obstruction / no-go note records the scoped verdict \texttt{Not Derived On Current Line}. That verdict does not weaken the benchmark package already in hand. It states only that the present frozen bridge object does not derive that package on its current assumptions.

The routing consequence is correspondingly narrow and explicit. Another deeper same-output relay below the same frozen bridge object no longer counts as the honest next benchmark-facing move. If derivational work is resumed, it must begin on a genuinely new line with an explicit observer-localizing law or equivalent new structural principle that says how substrate structure maps to observer spacetime fields, how the operative metric is generated locally, and how the Einsteinian observer equation together with relativistic recovery arise from that same local structure.

\section{Related Work and Literature Positioning}

A public-facing benchmark package should not leave its literature category to guesswork. The clean comparison is therefore explicit.
\begin{enumerate}
    \item \textbf{Standard GR and its congruence / ADM baselines.} Relative to ordinary GR and its classical congruence and Hamiltonian baselines \cite{Einstein1916,Raychaudhuri1955,EllisVanElst1999,ADM1962}, the active benchmark package keeps the same low-energy equations and observable content. The difference is not a second field equation, a new low-energy graviton sector, or a modified matter route. It is the disciplined interpretive overlay supplied by the \Aether{} / \Aflow{} ontology and by the congruence-based geometric reading of that ontology. On the same logic, the consistency module is an audit of the adopted GR sector rather than a claim of new low-energy dynamics.
    \item \textbf{Einstein-aether theory and other preferred-structure alternatives.} Einstein-aether theory and related preferred-structure models introduce a dynamical unit timelike vector, preferred-frame sector, or preferred foliation at observer level and therefore extra low-energy structure or mode content \cite{JacobsonMattingly2001,JacobsonMattingly2004,ElingJacobson2004,Jacobson2010,Jacobson2014,JacobsonSperanza2015,Horava2009}. The active benchmark package does not do that. It claims no surviving observer-level aether vector field, no extra low-energy graviton sector, and no modified Einsteinian observer law.
    \item \textbf{Emergent, induced, reconstruction, self-coupling, and analogue-gravity baselines.} Programs such as Jacobson's thermodynamic derivation, Sakharov-style induced gravity, canonical reconstruction, self-coupling consistency, Lovelock-type uniqueness analysis, and analogue-gravity comparison \cite{Jacobson1995,Visser2002,HojmanKucharTeitelboim1976,Deser2004,Lovelock1971,BarceloLiberatiVisser2011,Sakharov1967} are closer in ambition to the downstream derivation program than to the benchmark package itself. Their category is an attempt to explain why Einsteinian gravity appears, not an exact-GR interpretive package already separated from its open derivation burden.
\end{enumerate}

\begin{proposition}[Most accurate current literature label]
At current scope, \TheoryName{} is best described as an exact-GR interpretive package plus an open emergent-gravity research program, not as an already completed low-energy preferred-frame alternative to GR.
\end{proposition}

This proposition sharpens the only defensible novelty statement. The project does not earn novelty merely by saying that GR might emerge from deeper structure, because that ambition already belongs to an established emergent-gravity literature. It also does not earn novelty by introducing a new low-energy preferred-frame field theory, because the active benchmark package explicitly refuses that move. If a narrower novelty claim survives review, it lies instead in the combination of disciplined ontology, exact-GR benchmark packaging, benchmark-gated derivation architecture, and willingness to preserve a scoped no-go rather than relabel the failed frozen line as progress.

\begin{remark}
Representative comparison points therefore span three clusters: ordinary GR together with its congruence and ADM baselines, preferred-structure alternatives such as Einstein-aether and Ho\v{r}ava gravity, and the open emergent / induced / reconstruction conversation \cite{Einstein1916,Raychaudhuri1955,EllisVanElst1999,ADM1962,JacobsonMattingly2001,JacobsonMattingly2004,ElingJacobson2004,Jacobson2010,Jacobson2014,JacobsonSperanza2015,Horava2009,Jacobson1995,Visser2002,HojmanKucharTeitelboim1976,Deser2004,Lovelock1971,BarceloLiberatiVisser2011,Sakharov1967}.
\end{remark}

\section{The Flagship Exact-Closure Package}

The active public package should be read in the following fixed order.

\begin{center}
\begin{tabular}{p{0.07\textwidth} p{0.31\textwidth} p{0.52\textwidth}}
\toprule
\textbf{Order} & \textbf{Manuscript} & \textbf{Role} \\
\midrule
1 & Exact-Closure Sequence Overview & Canonical front door fixing the package role, reading order, and claim boundary. \\
2 & Exact-Closure Note & Short standalone anchor stating the exact-closure theory in its briefest complete form. \\
3 & Foundations & Ontological vocabulary, framework intent, and the exact-closure relation to GR. \\
4 & Dynamics & Effective action, matter coupling, weak-field structure, and benchmark exact law. \\
5 & Consistency & Gauge structure, degree-of-freedom counting, and effective health conditions. \\
6 & Relativistic Recovery & Exact relation to GR and local relation to SR. \\
7 & Flow Geometry & Congruence-based geometric dictionary for reading \Aflow{} inside exact relativistic structure. \\
\bottomrule
\end{tabular}
\end{center}

\begin{corollary}[Public packaging rule]
The overview-first exact-closure sequence is the flagship deliverable of the repository and should be presented before any downstream derivational continuation.
\end{corollary}

This order is not arbitrary. It places the finished effective theory statement before the open foundational program. Public accessibility and scientific honesty both improve when that order is held fixed.
The present flagship public article is built on that ordered core as a release-facing synthesis. It is not an eighth scientific source that reorders the package, and it does not displace the overview as the canonical front door.
The current-line obstruction / no-go note belongs downstream of this package in exactly that sense: it is a bounded verdict on a frozen bridge object, not a replacement front door and not a reason to blur the benchmark package back into open derivational work.

\section{Relation to the Downstream Derivational Bridge Archive}

The repository also contains a large derivational bridge archive. That archive remains important. It documents genuine attempts to recover or constrain the deeper substrate route, records negative and partial results, and preserves the history of the broader research program. Yet its role in public presentation is now different from what it would have been at an earlier stage.

The bridge archive is no longer the condition for saying that the repository has a theory statement. The theory statement is already given by \TheoryName{} as the exact-closure package. The archive is therefore downstream and optional. It should be read as foundational continuation answerable to the exact-closure benchmark, not as the front door to the project.

\begin{remark}
If a future bridge route succeeds, it must recover the same exact relativistic benchmark already fixed by the flagship package. If a future bridge route fails, the project contracts back to the exact-closure theory rather than back to ontology alone.
\end{remark}

\section{Research Provenance and Collaboration Method}

The repository should also be read honestly as AI-assisted theoretical research under sustained human direction. The original conceptual direction came from the user. GPT-5.4 assisted with drafting, mathematical exploration, branch screening, organization, and revision. The resulting package is therefore neither a purely conventional solo draft nor an autonomous AI production.

This provenance matters for transparency, but it does not change the intellectual standard to which the work should be held. The manuscripts should be judged by their assumptions, equations, internal logic, and honesty about what remains open. AI assistance does not convert speculative steps into established physics, and it does not remove the need for expert human review. Its role here has been methodological support within an explicitly documented collaboration.

\section{How the Package Should Be Reviewed}

At public-release stage, the most useful technical review focuses on five questions.
\begin{enumerate}
    \item Is the claim boundary stable and explicit throughout the package?
    \item Do the equations and prose say exactly what is established, and no more?
    \item Are the package roles of overview, anchor, and companion manuscripts clear and non-overlapping?
    \item Does the presentation distinguish properly between the flagship exact-closure package and the downstream bridge archive?
    \item Does the presentation make equally clear that the current bounded primitive-reservoir result is a scoped no-go on a frozen bridge object, so any resumed derivational work would require a genuinely new observer-localizing law rather than another same-output relay?
\end{enumerate}

Readers who need the shortest technically disciplined path should begin with the overview and then the exact-closure note. Readers who want the modular full statement should continue through foundations, dynamics, consistency, relativistic recovery, and flow geometry in that order. Only after that should the bridge archive be used for broader foundational context.

\section{Conclusion}

The repository now has a public-facing flagship theory package. The positive result is that \TheoryName{} can already be presented as a disciplined exact relativistic theory statement built on the \Aether{} / \Aflow{} ontology, with the effective gravitational dynamics adopted exactly as those of GR. The scope of that result is the overview-first exact-closure package and its companion manuscripts, not a completed first-principles substrate derivation.

The next honest burden is therefore twofold. First, public release materials should continue to package, review, and present this exact-closure package as the flagship result already achieved, with the present article serving as a release-facing synthesis rather than a reordering of the scientific core, and with the overview-first order kept explicit and easy to audit. Second, any future foundational continuation should remain clearly subordinate to that package, should state without ambiguity that deeper derivation from substrate variables remains open work rather than an achieved public claim, and should begin only if a genuinely new observer-localizing law or equivalent new structural principle is ready to be stated.

\begin{thebibliography}{99}

\bibitem{Einstein1916}
A.~Einstein,
``The Foundation of the General Theory of Relativity,''
\emph{Annalen der Physik} \textbf{49} (1916), 769--822.

\bibitem{ADM1962}
R.~Arnowitt, S.~Deser, and C.~W.~Misner,
``The Dynamics of General Relativity,''
in \emph{Gravitation: An Introduction to Current Research},
edited by L.~Witten (Wiley, New York, 1962), 227--265;
reprinted in \emph{General Relativity and Gravitation} \textbf{40} (2008), 1997--2027.

\bibitem{Raychaudhuri1955}
A.~Raychaudhuri,
``Relativistic Cosmology. I,''
\emph{Physical Review} \textbf{98} (1955), 1123--1126.

\bibitem{EllisVanElst1999}
G.~F.~R.~Ellis and H.~van Elst,
``Cosmological Models (Carg\`ese Lectures 1998),''
in \emph{Theoretical and Observational Cosmology},
edited by M.~Lachi\`eze-Rey,
NATO Science Series C: Mathematical and Physical Sciences Vol.~541
(Kluwer, Dordrecht, 1999), 1--116,
\href{https://arxiv.org/abs/gr-qc/9812046}{arXiv:gr-qc/9812046}.

\bibitem{JacobsonMattingly2001}
T.~Jacobson and D.~Mattingly,
``Gravity with a Dynamical Preferred Frame,''
\emph{Physical Review D} \textbf{64} (2001), 024028.

\bibitem{JacobsonMattingly2004}
T.~Jacobson and D.~Mattingly,
``Einstein-Aether Waves,''
\emph{Physical Review D} \textbf{70} (2004), 024003.

\bibitem{ElingJacobson2004}
C.~Eling and T.~Jacobson,
``Static Post-Newtonian Equivalence of General Relativity and Gravity with a Dynamical Preferred Frame,''
\emph{Physical Review D} \textbf{69} (2004), 064005.

\bibitem{Jacobson2010}
T.~Jacobson,
``Extended Ho\v{r}ava Gravity and Einstein-Aether Theory,''
\emph{Physical Review D} \textbf{81} (2010), 101502(R).

\bibitem{Jacobson2014}
T.~Jacobson,
``Undoing the Twist: The Ho\v{r}ava Limit of Einstein-Aether Theory,''
\emph{Physical Review D} \textbf{89} (2014), 081501(R).

\bibitem{JacobsonSperanza2015}
T.~Jacobson and A.~J.~Speranza,
``Variations on an Aethereal Theme,''
\emph{Physical Review D} \textbf{92} (2015), 044030.

\bibitem{Horava2009}
P.~Ho\v{r}ava,
``Quantum Gravity at a Lifshitz Point,''
\emph{Physical Review D} \textbf{79} (2009), 084008.

\bibitem{Jacobson1995}
T.~Jacobson,
``Thermodynamics of Spacetime: The Einstein Equation of State,''
\emph{Physical Review Letters} \textbf{75} (1995), 1260--1263.

\bibitem{Visser2002}
M.~Visser,
``Sakharov's Induced Gravity: A Modern Perspective,''
\emph{Modern Physics Letters A} \textbf{17} (2002), 977--992,
\href{https://arxiv.org/abs/gr-qc/0204062}{arXiv:gr-qc/0204062}.

\bibitem{HojmanKucharTeitelboim1976}
S.~A.~Hojman, K.~Kucha\v{r}, and C.~Teitelboim,
``Geometrodynamics Regained,''
\emph{Annals of Physics} \textbf{96} (1976), 88--135.

\bibitem{Deser2004}
S.~Deser,
``Self-Interaction and Gauge Invariance,''
\emph{General Relativity and Gravitation} \textbf{1} (1970), 9--18;
reprinted with 2004 note at
\href{https://arxiv.org/abs/gr-qc/0411023}{arXiv:gr-qc/0411023}.

\bibitem{Lovelock1971}
D.~Lovelock,
``The Einstein Tensor and Its Generalizations,''
\emph{Journal of Mathematical Physics} \textbf{12} (1971), 498--501.

\bibitem{BarceloLiberatiVisser2011}
C.~Barcelo, S.~Liberati, and M.~Visser,
``Analogue Gravity,''
\emph{Living Reviews in Relativity} \textbf{14} (2011), 3.

\bibitem{Sakharov1967}
A.~D.~Sakharov,
``Vacuum Quantum Fluctuations in Curved Space and the Theory of Gravitation,''
\emph{Soviet Physics Doklady} \textbf{12} (1968), 1040--1041
[translated from \emph{Doklady Akademii Nauk SSSR} \textbf{177} (1967), 70--71].

\end{thebibliography}


\end{document}
